% Options for packages loaded elsewhere
\PassOptionsToPackage{unicode}{hyperref}
\PassOptionsToPackage{hyphens}{url}
%
\documentclass[
  10pt,
]{article}
\usepackage{amsmath,amssymb}
\usepackage{iftex}
\ifPDFTeX
  \usepackage[T1]{fontenc}
  \usepackage[utf8]{inputenc}
  \usepackage{textcomp} % provide euro and other symbols
\else % if luatex or xetex
  \usepackage{unicode-math} % this also loads fontspec
  \defaultfontfeatures{Scale=MatchLowercase}
  \defaultfontfeatures[\rmfamily]{Ligatures=TeX,Scale=1}
\fi
\usepackage{lmodern}
\ifPDFTeX\else
  % xetex/luatex font selection
\fi
% Use upquote if available, for straight quotes in verbatim environments
\IfFileExists{upquote.sty}{\usepackage{upquote}}{}
\IfFileExists{microtype.sty}{% use microtype if available
  \usepackage[]{microtype}
  \UseMicrotypeSet[protrusion]{basicmath} % disable protrusion for tt fonts
}{}
\makeatletter
\@ifundefined{KOMAClassName}{% if non-KOMA class
  \IfFileExists{parskip.sty}{%
    \usepackage{parskip}
  }{% else
    \setlength{\parindent}{0pt}
    \setlength{\parskip}{6pt plus 2pt minus 1pt}}
}{% if KOMA class
  \KOMAoptions{parskip=half}}
\makeatother
\usepackage{xcolor}
\usepackage[margin=22mm]{geometry}
\usepackage{longtable,booktabs,array}
\usepackage{calc} % for calculating minipage widths
% Correct order of tables after \paragraph or \subparagraph
\usepackage{etoolbox}
\makeatletter
\patchcmd\longtable{\par}{\if@noskipsec\mbox{}\fi\par}{}{}
\makeatother
% Allow footnotes in longtable head/foot
\IfFileExists{footnotehyper.sty}{\usepackage{footnotehyper}}{\usepackage{footnote}}
\makesavenoteenv{longtable}
\setlength{\emergencystretch}{3em} % prevent overfull lines
\providecommand{\tightlist}{%
  \setlength{\itemsep}{0pt}\setlength{\parskip}{0pt}}
\setcounter{secnumdepth}{-\maxdimen} % remove section numbering
\ifLuaTeX
\usepackage[bidi=basic]{babel}
\else
\usepackage[bidi=default]{babel}
\fi
\babelprovide[main,import]{italian}
% get rid of language-specific shorthands (see #6817):
\let\LanguageShortHands\languageshorthands
\def\languageshorthands#1{}
\ifLuaTeX
  \usepackage{selnolig}  % disable illegal ligatures
\fi
\usepackage{bookmark}
\IfFileExists{xurl.sty}{\usepackage{xurl}}{} % add URL line breaks if available
\urlstyle{same}
\hypersetup{
  pdflang={it},
  hidelinks,
  pdfcreator={LaTeX via pandoc}}

\author{}
\date{}

\begin{document}

\section{DDTA - Working Metamodel - Macro
Requirement}\label{ddta---working-metamodel---macro-requirement}

\begin{quote}
\textbf{NON CANONICAL WORKING DRAFT}

This file defines the general document metamodel. ThreatForge-specific
observations are excluded and belong to the case-study material.
\end{quote}

\subsection{1. Scope}\label{1-scope}

The metamodel describes the knowledge of a governed project before
deciding its concrete Doc-as-Code serialization and before deciding how
a software engine validates or assists authoring.

The governed product remains documentation-as-code: versionable,
diffable, reviewable files (for example YAML registries and Markdown
bodies). A software tool may apply and verify the metamodel, but it is
not the subject described by the project\textquotesingle s documents.

\subsubsection{Layer separation}\label{layer-separation}

\begin{enumerate}
\def\labelenumi{\arabic{enumi}.}
\tightlist
\item
  \textbf{Domain metamodel} - meaning of documents, entities and
  relations.
\item
  \textbf{Doc-as-Code representation} - how that meaning is serialized
  in files.
\item
  \textbf{Engine} - how tools load, validate, correlate, project and
  assist those files.
\end{enumerate}

\subsection{2. Principles already
fixed}\label{2-principles-already-fixed}

\subsubsection{Principle 1 - Progressive specialization and semantic
containment}\label{principle-1---progressive-specialization-and-semantic-containment}

Governed project knowledge is represented as a hierarchy of
progressively specialized governed documents. A Macro Requirement
establishes a macro-level subject and purpose; Decisions narrow that
subject through explicit choices; Functional Requirements translate
those choices into independently verifiable obligations; Specialized
Requirements add domain-specific constraints to one Functional
Requirement.

Hierarchical parentage represents semantic containment and remains
distinct from non-hierarchical traceability, evidence and provenance
relations.

\subsubsection{Principle 2 - Identity independent from textual
encoding}\label{principle-2---identity-independent-from-textual-encoding}

Every governed entity has a stable identity independent from the textual
encoding of that identity. Hierarchy is established by explicit semantic
relationships. Composite identifier formats may mirror ancestry but must
not be the authoritative source of topology.

\subsection{3. Progressive disclosure by
competence}\label{3-progressive-disclosure-by-competence}

The hierarchy increases both specificity and the amount of vertical
competence required from the reader.

\begin{itemize}
\tightlist
\item
  \textbf{Macro Requirement} - understandable by stakeholders who
  understand the project domain, including non-software stakeholders.
\item
  \textbf{Decision} - requires more knowledge of the solution space and
  project choices.
\item
  \textbf{Functional Requirement} - precise, testable project behaviour
  or outcome.
\item
  \textbf{Specialized Requirement} - vertical competence such as
  governance, security, quality or another project-specific discipline.
\item
  \textbf{Analysis models and methodology-specific records} - specialist
  analytical concepts introduced only when needed.
\end{itemize}

\subsubsection{Executive-map criterion}\label{executive-map-criterion}

Reading only the \textbf{title and Intent} of all Macro Requirements of
a governed project must produce a recognizable summary of:

\begin{itemize}
\tightlist
\item
  what the project is trying to realize;
\item
  what general value it produces;
\item
  why its main blocks of work exist.
\end{itemize}

This criterion applies to the \textbf{governed project}, not to the tool
used to govern it.

\subsection{4. Domain vocabulary versus metamodel
vocabulary}\label{4-domain-vocabulary-versus-metamodel-vocabulary}

The metamodel must not ban technical words merely because they are
technical.

A term belongs in an MR when it is part of the \textbf{application
domain} and is reasonably understandable to the stakeholders of that
domain. For example, a cryptography project may naturally use security
terminology; a medical imaging project may naturally use imaging
terminology; a facial-recognition access project may naturally use terms
such as facial recognition, biometric data or ML model.

What the MR must not be forced to contain are technical concepts
introduced \textbf{only because the documentation model or an analysis
method expects them}, such as:

\begin{itemize}
\tightlist
\item
  component inventories;
\item
  data-flow decompositions;
\item
  trust boundaries;
\item
  protocol details;
\item
  attack-surface classifications;
\item
  STRIDE/LINDDUN categories;
\item
  method-specific threat or privacy taxonomies.
\end{itemize}

Those concepts normally emerge later, when Decisions, Requirements, a
project/system representation or methodology-specific analysis require
them.

\subsection{5. Macro Requirement
definition}\label{5-macro-requirement-definition}

A Macro Requirement identifies one major project concern or
macro-capability. It explains why that concern matters, who is affected,
and the general result the project seeks at that level.

It is the root document type of the governed-document hierarchy.

\subsubsection{Fundamental question}\label{fundamental-question}

\begin{quote}
\textbf{Which major project concern or result are we governing, why does
it matter, and who is involved?}
\end{quote}

\subsection{6. Candidate information
model}\label{6-candidate-information-model}

\begin{longtable}[]{@{}lll@{}}
\toprule\noalign{}
Concept & State & Purpose \\
\midrule\noalign{}
\endhead
\bottomrule\noalign{}
\endlastfoot
\texttt{id} & required & stable identity, history and traceability \\
\texttt{title} & required & concise name understandable within the
project domain \\
\texttt{lifecycle} & required & current/historical state of the
document \\
\texttt{intent} & required & concise statement of the macro purpose,
value and desired outcome that justify the MR \\
\texttt{context} & required & minimum background needed to understand
the concern \\
\texttt{stakeholders} & required & people/groups that benefit, use,
govern or are affected \\
\texttt{scope} & open semantic decision & boundary of the macro concern;
exact form still under study \\
\texttt{assumptions} & optional & facts accepted as true that materially
affect the whole MR \\
\texttt{constraints} & optional & limits that materially constrain the
whole MR \\
\texttt{non\_goals} & open semantic decision & plausible expectations
deliberately outside the MR, if useful \\
\end{longtable}

\subsubsection{Concision}\label{concision}

A Macro Requirement must be effective rather than verbose. It should not
become a checklist of implementation rules.

\begin{itemize}
\tightlist
\item
  \texttt{title}: short noun phrase or short descriptive title;
\item
  \texttt{intent}: normally 1-2 sentences covering macro purpose, value
  and desired outcome;
\item
  \texttt{context}: only the background needed to interpret the MR;
\item
  \texttt{stakeholders}: macro stakeholder roles, not a technical
  inventory;
\item
  assumptions/constraints: only if they affect the whole branch below
  the MR.
\end{itemize}

\subsection{7. Decisions already
closed}\label{7-decisions-already-closed}

\subsubsection{7.1 Intent owns the macro purpose, value and desired
outcome}\label{71-intent-owns-the-macro-purpose-value-and-desired-outcome}

\texttt{Objective} is \textbf{removed} from the Macro Requirement
conceptual model. Across the four facial-access examples and the five
ThreatForge case-study MRs, a separate Objective did not show stable
autonomous semantics: it either paraphrased the Intent or began
decomposing it into Decision/Requirement-level behavior.

\texttt{Intent} therefore answers both \textbf{why the macro concern
exists} and \textbf{what macro outcome/value it seeks}. It remains
concise and must not become a checklist of atomic obligations.

Use the \textbf{Intent-completeness test}: if \texttt{title\ +\ Intent}
are insufficient for a domain-competent stakeholder to understand the
macro concern and desired result, improve the Intent rather than adding
a second Objective field.

Use the \textbf{Objective-leakage test} when reviewing legacy or draft
material: if a separate objective merely restates Intent, merge the
useful wording into Intent; if it enumerates specific behaviors,
mechanisms or independently verifiable outcomes, move that content to
Scope, Constraints, Decisions or Requirements as appropriate.

\subsubsection{7.2 Stakeholders are semantically
required}\label{72-stakeholders-are-semantically-required}

The MR must identify the stakeholder groups relevant to its macro
concern. Their names may be technical when that is normal domain
language, but they should be meaningful to the project community rather
than internal implementation labels.

\subsubsection{7.3 Assumptions and Constraints are
distinct}\label{73-assumptions-and-constraints-are-distinct}

\begin{itemize}
\tightlist
\item
  \textbf{Assumption}: something currently treated as true for the
  branch of work.
\item
  \textbf{Constraint}: a limit within which the branch must operate.
\end{itemize}

Both are optional and belong in the MR only when they materially
influence the whole branch below it.

\subsubsection{7.4 Parent and children}\label{74-parent-and-children}

\texttt{MacroRequirement.parent\ =\ null}.

A Macro Requirement may own \texttt{0..*} Decisions. The child
collection is derived from
\texttt{Decision.parent\ -\textgreater{}\ MacroRequirement}; it is not a
second topology authority.

\subsubsection{7.5 Cross-MR dependency}\label{75-cross-mr-dependency}

A Macro Requirement may have a non-hierarchical
\texttt{dependsOn\ -\textgreater{}\ MacroRequirement\ {[}0..*{]}}
relation when another macro concern is a genuine dependency. The exact
Doc-as-Code serialization is deferred.

\subsubsection{7.6 No Acceptance Criteria at MR
level}\label{76-no-acceptance-criteria-at-mr-level}

The MR does not contain atomic acceptance criteria. Evidence of
satisfaction emerges from its descendant Requirements and their
verification evidence.

\subsubsection{7.7 Progressive
technicality}\label{77-progressive-technicality}

The MR should be understandable by stakeholders competent in the
\textbf{project domain} without requiring them to understand
implementation decomposition or analysis methodology. More specialized
technical knowledge is introduced as documentation descends.

\subsubsection{7.8 Technical and analytical decomposition normally
emerges
later}\label{78-technical-and-analytical-decomposition-normally-emerges-later}

The MR must not be structurally required to enumerate components,
services, APIs, data flows, trust boundaries, attack surfaces or threat
categories. Those normally emerge at lower documentation or analysis
levels.

This does \textbf{not} prohibit domain-specific terminology in an MR.

\subsubsection{7.9 MR prose must be stable with respect to design
progress}\label{79-mr-prose-must-be-stable-with-respect-to-design-progress}

A Macro Requirement describes a durable project concern, not the current
progress of solution design. Statements such as
\texttt{not\ yet\ decided}, \texttt{currently\ evaluating},
\texttt{will\ be\ chosen\ later} or equivalent design-status prose
normally belong to working notes, Decision lifecycle/status or another
planning artifact rather than to the MR body.

Use the \textbf{temporal-stability test}: if a sentence becomes obsolete
merely because a lower-level Decision is taken, while the macro project
concern remains unchanged, the sentence should normally not be part of
the MR.

This rule does not forbid time-dependent \textbf{domain facts}. For
example, an authorization may expire or a biometric enrollment may have
a lifecycle because those are properties of the application domain. The
rule targets temporary authoring/design state, not legitimate temporal
semantics of the project.

\subsubsection{7.10 Teachability and macro project-map
projection}\label{710-teachability-and-macro-project-map-projection}

A Macro Requirement must contribute to a project description that can be
\textbf{taught progressively}, not merely validated. Individually, each
MR should provide one understandable conceptual block through its
identity, title and Intent. Collectively, the MR set should support a
simple didactic map of the project\textquotesingle s major concerns and
their meaningful relationships without requiring lower-level
architecture knowledge.

Use the \textbf{teachability test}: a reader who understands the
application domain but is new to the project should be able to use the
MR set to build a simple mental/graphical map answering:

\begin{enumerate}
\def\labelenumi{\arabic{enumi}.}
\tightlist
\item
  what the main project concerns are;
\item
  why each concern exists;
\item
  which macro concerns depend on, support or otherwise relate to others
  when such relations are explicitly governed or can be shown as clearly
  marked non-canonical teaching interpretations;
\item
  how the blocks collectively explain the project\textquotesingle s
  macro functioning without exposing lower-level implementation
  topology.
\end{enumerate}

Failure of this test is evidence to review MR decomposition, titles,
Intent wording or missing explicit cross-MR relations. It is not
permission to invent architecture in order to make the picture easier to
draw.

\subsubsection{7.11 Analysis-enabling does not mean
analysis-producing}\label{711-analysis-enabling-does-not-mean-analysis-producing}

MR content can provide context and documentary evidence useful to later
system representation and analysis, but an MR does \textbf{not}
instantiate canonical analytical elements merely because it names a
person, device, datum, ML capability or other domain concept.

A didactic map derived from MRs is therefore not a DFD, architecture
diagram or threat-analysis model. Canonical components, actors, data
resources, data flows, trust boundaries or method-specific
interpretations require later governed modeling/review.

\subsection{8. MR invariants}\label{8-mr-invariants}

\begin{enumerate}
\def\labelenumi{\arabic{enumi}.}
\tightlist
\item
  \textbf{Root hierarchy} - MR has no parent.
\item
  \textbf{Stable identity} - identity is distinct from its textual
  encoding.
\item
  \textbf{Single macro concern} - the MR expresses one coherent macro
  responsibility.
\item
  \textbf{Child containment} - descendants progressively narrow the MR
  rather than introduce unrelated scope.
\item
  \textbf{No implementation choice} - choices of solution belong to
  Decision unless they are intrinsic to the project domain itself.
\item
  \textbf{No requirement-level obligation} - atomic testable obligations
  belong to Requirements.
\item
  \textbf{Broad audience within the domain} - title and Intent can be
  understood by non-software stakeholders who understand the project
  domain.
\item
  \textbf{Project-map usefulness} - titles + Intent of all MR summarize
  the project and its major work blocks.
\item
  \textbf{Explicit dependencies} - cross-MR dependencies are not hidden
  as pseudo-hierarchy.
\item
  \textbf{Analysis enabling, not analysis producing} - MR provides
  useful project context without directly instantiating canonical
  analytical elements or prescribing an analysis paradigm.
\item
  \textbf{Architecture resilience} - changing a lower-level architecture
  choice should not normally force the MR to be rewritten if the macro
  project concern has not changed.
\item
  \textbf{Temporal stability} - taking or revising a lower-level
  Decision should not make MR prose obsolete unless the macro project
  concern itself changes.
\item
  \textbf{Teachability} - each MR is an understandable macro block and
  the MR set can support a simple didactic project map without requiring
  lower-level architecture.
\item
  \textbf{Projection honesty} - a teaching projection derived from MRs
  must not be presented as a DFD, architecture model or canonical
  analysis representation.
\end{enumerate}

\subsection{9. Semantic decisions still
open}\label{9-semantic-decisions-still-open}

The Intent/Objective question is now closed: \texttt{Objective} is
removed and \texttt{Intent} owns macro purpose, value and desired
outcome. The remaining MR questions below must still be tested before
final freeze.

\subsubsection{9.1 Meaning and shape of
Scope}\label{91-meaning-and-shape-of-scope}

We need examples to determine what \texttt{scope} should actually
express and whether \texttt{included}, \texttt{excluded} and
\texttt{non\_goals} are all necessary or partly redundant.

The goal is not to create three places that repeat the same boundary in
different words.

\subsubsection{9.2 When to split or merge Macro
Requirements}\label{92-when-to-split-or-merge-macro-requirements}

MR boundaries must not be decided from wording similarity alone. This is
especially important when refactoring an existing governed corpus: two
rewritten MRs can sound similar because both were described through the
same downstream consumer, while their descendant Decisions still govern
different project concerns. Conversely, historical separation is not
sufficient evidence to keep two MRs if it only mirrors an implementation
split.

Before splitting or merging candidate MRs, test the boundary with these
questions:

\begin{enumerate}
\def\labelenumi{\arabic{enumi}.}
\tightlist
\item
  \textbf{Independent value} - does each candidate express a
  project-level result that remains meaningful without describing the
  other?
\item
  \textbf{Decision-family coherence} - where an existing corpus already
  has Decisions, do the descendants of each candidate form distinct and
  internally coherent families of choices?
\item
  \textbf{Solution-vocabulary removal} - does the distinction survive
  after removing names of current components, schemas, tools, analysis
  representations and other chosen mechanisms?
\item
  \textbf{Dependency versus containment} - can the relation be expressed
  more accurately as \texttt{dependsOn}, rather than by merging concerns
  or introducing false hierarchy?
\item
  \textbf{Stakeholder explanation} - can a stakeholder competent in the
  project domain explain the difference using title + Intent without
  knowing the implementation?
\item
  \textbf{Merge stress test} - would a merged MR need to explain
  substantially different families of Decisions or different kinds of
  project value?
\item
  \textbf{Split stress test} - is the separation justified by project
  concerns, rather than merely by architectural layers, subsystems or
  workflow stages?
\end{enumerate}

These are currently \textbf{boundary-test heuristics}, not yet frozen
deterministic invariants. They must be exercised on multiple project
examples.

\paragraph{Refactoring rule for existing
documentation}\label{refactoring-rule-for-existing-documentation}

When an MR already owns a body of Decisions, reconstruct its macro
concern from the semantic responsibility of that Decision family before
rewriting the MR. A rewritten MR should explain why those Decisions
belong together without merely summarizing their current solution
vocabulary.

\subsection{10. Derived didactic projection - Macro Project
Map}\label{10-derived-didactic-projection---macro-project-map}

The MR set may support a \textbf{Macro Project Map}: a non-canonical,
teaching-oriented projection that gives readers a block-level view of
the project\textquotesingle s major concerns. Its purpose is orientation
and progressive disclosure, not system analysis.

A candidate projection uses:

\begin{itemize}
\tightlist
\item
  one block per MR, showing canonical MR identity and title and
  optionally a shortened Intent;
\item
  explicit cross-MR relations such as \texttt{dependsOn} when available;
\item
  clearly labeled didactic links only when they are visibly marked as
  non-canonical interpretation and traceable to supporting MR text;
\item
  edge labels expressed as verbs/relations rather than unlabeled arrows;
\item
  textual equivalent/traceability alongside the graphic.
\end{itemize}

The projection must \textbf{not} reinterpret MR blocks as DFD processes,
external entities or data stores, and its links are not data flows.
Layout, color, coordinates and SVG/HTML are renderer concerns and must
remain outside the semantic projection.

\subsubsection{Accessibility / BES-oriented presentation
heuristic}\label{accessibility--bes-oriented-presentation-heuristic}

The teaching view should favor low cognitive load and progressive
disclosure:

\begin{enumerate}
\def\labelenumi{\arabic{enumi}.}
\tightlist
\item
  few blocks per view;
\item
  canonical titles, not unexplained synonyms;
\item
  one dominant idea per block;
\item
  every arrow labeled in text;
\item
  meaning never encoded by color alone;
\item
  a textual explanation equivalent to the visual map;
\item
  stable navigation from a block back to its MR.
\end{enumerate}

\subsubsection{Rendering compatibility observation from
ThreatForge}\label{rendering-compatibility-observation-from-threatforge}

At ThreatForge planning baseline
\texttt{44f77796469465a6a076596b1b1ba1afdde92db7}, the active Base DFD
path already follows a useful pattern: renderer-neutral semantic
projection, deterministic layout, and self-contained HTML with embedded
accessible SVG and traceability metadata. The active implementation is a
custom Node/ESM renderer rather than Graphviz or Mermaid.

This study therefore records a \textbf{future reuse direction}, not an
implementation decision: a Macro Project Map renderer may share visual
tokens, accessibility patterns, deterministic layout utilities and
HTML/SVG shell with the DFD renderer while keeping a separate
MR-specific semantic contract.

\subsection{11. Observation deferred to the future analysis
metamodel}\label{11-observation-deferred-to-the-future-analysis-metamodel}

Concrete MR examples show that different concerns inside the same
governed project may later benefit from different analysis paradigms. An
ML-based capability may justify AI/ML-oriented threat analysis, while
another concern in the same project may justify privacy-oriented
analysis or a different specialist method.

This observation does \textbf{not} change the MR contract and must not
introduce method-specific fields into Macro Requirements. The project
hierarchy describes project concerns; analysis methodology is a later
concern.

The future analysis metamodel should therefore test whether it can:

\begin{enumerate}
\def\labelenumi{\arabic{enumi}.}
\tightlist
\item
  select an analysis method for an explicit analysis scope or set of
  governed subjects rather than assigning one methodology to an MR by
  construction;
\item
  allow different methods to analyze different, overlapping or
  complementary parts of one project;
\item
  preserve method-specific interpretation without contaminating project
  knowledge with methodology vocabulary;
\item
  converge findings, documentation gaps, proposed requirements and other
  feedback into one governed project/documentation lifecycle.
\end{enumerate}

These are \textbf{deferred design questions}, not frozen analysis-model
decisions. They are recorded here so that the MR model does not
accidentally preclude multi-method project analysis.

\subsection{12. Doc-as-Code mapping deliberately
deferred}\label{12-doc-as-code-mapping-deliberately-deferred}

The following remain representation questions, not domain-model
decisions:

\begin{itemize}
\tightlist
\item
  where each MR concept lives between registry and Markdown body;
\item
  whether \texttt{parent\ =\ null} is serialized or implicit;
\item
  concrete representation of stakeholders, assumptions and constraints;
\item
  concrete representation of \texttt{dependsOn};
\item
  final Markdown heading names;
\item
  deterministic checks that can enforce the semantic contract without
  pretending to validate human meaning.
\end{itemize}

\textbf{STOP:} Decision is not modeled in this document yet.

\end{document}
